% !TeX root = ../navier-stokes-manuscript.tex
% ============================================================
% 2.6 Finite Adjacent Rectangles
% ============================================================

\subsection{Finite Adjacent Rectangles}\label{sub:FiniteAdjacentRectangles}

The whole intersection is reached through finite estimates, so begin with the first finite demand that cannot be met by one isolated Sobolev membership.
Fix \(I=(0,T)\) with \(T<T_*\), and ask that both
\[
  V_0=u
  \qquad\text{and}\qquad
  V_1=\partial_tu
\]
arrive at the same time slice in \(H^2(\R^3)\).
An \(L^2_tH^3_x\) bound for \(V_0\) supplies spatial strength but no time trace by itself.
The missing temporal side is
\[
  \partial_tV_0=V_1\in L^2(I;H^1).
\]
To place \(V_1\) at the same \(H^2\) trace height, the same argument asks for
\[
  V_1\in L^2(I;H^3),
  \qquad
  \partial_tV_1=V_2\in L^2(I;H^1).
\]
The four requirements are
\[
  V_0,V_1\in L^2(I;H^3),
  \qquad
  V_1,V_2\in L^2(I;H^1).
\]
The repeated \(V_1\) is not redundant: it is the derivative of \(V_0\) in the first trace pair and the field whose own derivative is \(V_2\) in the second.

The trace mechanism explains why the two spatial levels must be adjacent around \(H^2\).
Let
\[
  A:=(1-\Delta)^{1/2},
  \qquad
  \norm{v}_{H^s}=\norm{A^sv}_2.
\]
For a smooth \(v\),
\[
  \norm{v(t)}_{H^2}^2-\norm{v(s)}_{H^2}^2
  =
  2\operatorname{Re}
  \int_s^t
  \left\langle
    A\partial_\tau v(\tau),
    A^3v(\tau)
  \right\rangle_{L^2}
  \,d\tau.
\]
The integrand satisfies
\[
  \left|
    \left\langle
      A\partial_\tau v,
      A^3v
    \right\rangle_{L^2}
  \right|
  \le
  \norm{\partial_\tau v}_{H^1}
  \norm{v}_{H^3},
\]
which belongs to \(L^1(I)\) by Cauchy--Schwarz when
\[
  v\in L^2(I;H^3),
  \qquad
  \partial_tv\in L^2(I;H^1).
\]
The standard density argument then yields
\[
  v\in C\bigl([0,T];H^2(\R^3)\bigr).
\]
Applied first to \(V_0\) and then to \(V_1\), this gives
\[
  V_0,V_1
  \in
  C\bigl([0,T];H^2(\R^3)\bigr).
\]
The two trace calculations use four norms, which we collect in
\begin{align*}
  \mathfrak R_{1,2}(u;T)
  :={}&
  \norm{u}_{L^2(0,T;H^3)}^2
  +
  \norm{\partial_tu}_{L^2(0,T;H^3)}^2
  \\
  &+
  \norm{\partial_tu}_{L^2(0,T;H^1)}^2
  +
  \norm{\partial_t^2u}_{L^2(0,T;H^1)}^2.
\end{align*}
Every entry is a derivative of the original solution.
The rectangle groups the four norms needed for the two traces; it does not generate new velocity fields.

The same trace calculation determines the general cut.
Fix independent integers \(N,M\in\Nzero\).
The desired trace height is \(H^M\), and its adjacent pair is
\[
  H^{M+1}(\R^3)
  \hookrightarrow
  H^M(\R^3)
  \hookrightarrow
  H^{M-1}(\R^3).
\]
When \(M=0\), the lower member is the inhomogeneous negative space \(H^{-1}(\R^3)\).
For every \(0\le n\le N\), the trace of \(V_n\) at height \(H^M\) asks for
\[
  V_n\in L^2(0,T;H^{M+1}),
  \qquad
  V_{n+1}=\partial_tV_n\in L^2(0,T;H^{M-1}).
\]
The finite velocity norm containing all of these adjacent pairs is
\begin{align*}
  \mathfrak R_{N,M}(u;T)
  :={}&
  \sum_{n=0}^{N}
  \norm{V_n}_{L^2(0,T;H^{M+1})}^{2}
  \\
  &+
  \sum_{n=0}^{N}
  \norm{V_{n+1}}_{L^2(0,T;H^{M-1})}^{2}.
\end{align*}
For a smooth representative, each retained temporal row obeys
\[
  \left|
    \frac{d}{dt}\norm{V_n(t)}_{H^M}^{2}
  \right|
  \le
  2
  \norm{V_{n+1}(t)}_{H^{M-1}}
  \norm{V_n(t)}_{H^{M+1}}.
\]
The right-hand side lies in \(L^1(0,T)\), and density together with weak continuity yields
\[
  V_n
  \in
  C\bigl([0,T];H^M(\R^3)\bigr),
  \qquad
  0\le n\le N.
\]
Increasing \(N\) retains more temporal responses.
Increasing \(M\) raises the spatial trace height.
The two choices remain independent, and the bound attached to the rectangle may depend on both.

The number \(\mathfrak R_{N,M}(u;T)\) measures the velocity rows needed for those traces.
To compare finite blocks without losing the equation that relates their entries, keep the corresponding pressure rows and differentiated equations as well.
For the solution \((u,p)\) generated by \(u_0\), define
\[
  \cR_{N,M}[u_0;T]
  :=
  \left(
    (V_n)_{n=0}^{N+1},
    (Q_n)_{n=0}^{N};
    \mathfrak R_{N,M}(u;T)
  \right),
\]
where
\[
  V_n=\partial_t^nu,
  \qquad
  Q_n=\partial_t^np.
\]
For every \(0\le n\le N\), the retained rows satisfy
\[
  V_{n+1}
  +
  \sum_{a=0}^{n}
  \binom{n}{a}
  (V_a\cdot\nabla)V_{n-a}
  +\nabla Q_n
  =
  \nu\Delta V_n,
\]
\[
  \nabla\cdot V_n=0,
\]
and
\[
  -\Delta Q_n
  =
  \partial_i\partial_j
  \sum_{a=0}^{n}
  \binom{n}{a}
  (V_a)_i(V_{n-a})_j.
\]
Rapid decay fixes the normalization of every \(Q_n\).
At the retained Sobolev depths, each mixed velocity coordinate is the weak spatial derivative
\[
  J_{n,\alpha}=\partial_x^\alpha V_n,
\]
and its pressure coordinate is the corresponding derivative of \(Q_n\).
Thus \(\mathfrak R_{N,M}\) is the finite velocity norm, while \(\cR_{N,M}\) also retains the equations and normalized pressure rows that identify those velocities as derivatives of \((u,p)\).
We will compare the same finite block on different time intervals.
For \(J=(0,\tau)\) with \(0<\tau\leq T_*\), let \(\mathscr X_{N,M}(J)\) denote the records whose velocity rows satisfy
\[
  W_n\in L^2(J;H^{M+1}),
  \qquad
  W_{n+1}\in L^2(J;H^{M-1})
  \qquad(0\leq n\leq N),
\]
with \(\partial_tW_n=W_{n+1}\) in distributions on \(J\times\R^3\) for \(0\leq n\leq N\),
whose normalized pressure rows satisfy
\[
  P_n
  =R_iR_j\sum_{a=0}^{n}\binom{n}{a}(W_a)_i(W_{n-a})_j
  \in L^1(J;H^M)
  \qquad(0\leq n\leq N),
\]
with \(\nabla P_n\in L^1(J;H^{M-1})\),
and whose retained rows obey the displayed momentum and divergence equations in distributions on \(J\times\R^3\).
The scalar component of such a record is the adjacent-row norm computed from its velocity rows on \(J\).
In particular,
\[
  \cR_{N,M}[u_0;T]
  \in
  \mathscr X_{N,M}((0,T)).
\]

As \((N,M)\) ranges over the finite pairs, these rectangles cover every finite coordinate of the velocity intersection.
Each rectangle has its own finite constant.
To assemble them, we have to verify that whenever two rectangles contain the same derivative, they contain the same distribution.
